\section{Finite Fourier coefficients in a logarithmic variable}
\label{sec:fourier-coefficients}

The polynomial coefficient algebra extends naturally to finite sums
\[
 B(L)=\sum_{\omega\in\Omega}P_\omega(L)e^{i\omega L},
 \qquad L=\log u,\qquad \Omega\subset\R\text{ finite},
\]
where every $P_\omega$ is a polynomial and the frequencies are fixed real
numbers. This includes $\sin(\log u)$, several incommensurable frequencies,
and polynomial amplitudes multiplying those oscillations. The source powers
and the Fourier frequencies have different roles: source weights determine
asymptotic truncation, while frequencies label coefficient modes. No Fourier
mode is discarded merely because its frequency is large.

\subsection{Coefficient arithmetic and reality}

Multiplication convolves frequencies, and the Euler derivative is
\[
 \frac{\dd}{\dd L}\bigl(P_\omega(L)e^{i\omega L}\bigr)
 =\bigl(P_\omega'(L)+i\omega P_\omega(L)\bigr)e^{i\omega L}.
\]
Thus this algebra is closed under precisely the operations used by
Theorem~\ref{thm:coefficients}. Equal frequencies are identified before
their polynomial amplitudes are added. Frequency equality is an exact
mathematical relation; an approximation of two close frequencies does not
justify treating them as equal.

The real-valued condition is the conjugacy relation
$P_{-\omega}(L)=\overline{P_\omega(L)}$ on real $L$. This relation is the condition imposed on any fixed
parameters in the amplitudes. Complex exponential notation is compatible
with a real-valued coefficient because conjugate modes combine into real
trigonometric expressions. Affine trigonometric phases, such as $\sin(2L+1)$, contribute their
constant phase factor to the polynomial amplitude.

\subsection{Inverse coefficients and independent composition}

The admitted forward family is
\[
 f(u)=y_0+a u^p\left(1+\sum_j u^{\delta_j}B_j(\log u)\right),
 \qquad ap\ne0,\quad \delta_j>0,
\]
with finitely many real Fourier-polynomial coefficients. The leading core is
a monomial with a proved nonzero real coefficient. An oscillatory leading
block, even one that happens to remain positive, requires a different inverse
normalization and is not covered by this construction.

For a multi-index $\bk$ of depth $n$ and weight
$\sigma=\bk\cdot\bdelta$, the unchanged coefficient formula is
\[
 C^{(r)}_{\bk}(L)=\frac{(-1)^n r}{p^n\bk!}
 \prod_{j=1}^{n-1}(\D+r+\sigma+pj)\prod_i B_i(L)^{k_i}.
\]
Convolution forms the product, and the Fourier Euler operation applies each
factor. Contributions at a common source weight must then be combined
to obtain a complete asymptotic block.

Residual checking uses a separate composition recurrence. If
$u=z(1+U)$, then
\[
 (1+U)^\alpha B(L+\log(1+U))=\sum_{k\ge0}Q_k(L)U^k,
 \qquad Q_0=B,\qquad
 Q_{k+1}=\frac{(\alpha-k)Q_k+\D Q_k}{k+1}.
\]
This follows by Taylor expansion in the multiplicative increment; it holds
for each exponential-polynomial mode as well as for their sum. Its use in
the normalized forward equation independently checks the inverse coefficients.

As in Lemma~\ref{lem:composition}, an identically zero complete coefficient
is absorbing: if $Q_{k+1}=0$, then $\D Q_{k+1}=0$, so the homogeneous
recurrence makes every later $Q_j$ zero. The identity must hold after equal
frequencies and their full polynomial amplitudes have been combined, under
the fixed parameter hypotheses; a zero at one value of $L$ is insufficient.
Independently, if $\nu=\val U>0$, every term of $U^j$ has source weight at
least $j\nu$. Hence $(k+1)\nu\ge H$ excludes all $Q_jU^j$ with $j\ge k+1$
from the expansion at exclusive source-weight cutoff $H$. This removes
their retained blocks, not their contribution to a remainder estimate:
an omitted block at weight $H$ still needs its logarithmic envelope.
An integer $\alpha$ alone does not imply termination. For $B=e^{iL}$ and
$\alpha=0$, one has $Q_k=\binom{i}{k}e^{iL}\ne0$ for every $k\ge0$.
Even a nonconstant zero-frequency amplitude can prevent termination:
for $B=L$ and $\alpha=1$, $Q_1=L+1$ and $Q_2=1/2$.

For a nonunit observable power, the residual first reconstructs the source
unit by raising the observable unit to its reciprocal power.

\subsection{Envelopes, convergence, and omitted terms}

For real $L$, Fourier modes have modulus one. If
$P_\omega(L)=\sum_j c_{\omega,j}L^j$, then
\[
 |B(L)|\le\sum_{\omega,j}|c_{\omega,j}|(1+|L|)^j.
\]
This gives a nonoscillatory envelope for every retained block. Euler
differentiation preserves the polynomial degree; frequency factors affect
coefficient constants, rather than the power of the logarithmic envelope.
Consequently a multi-index coefficient has polynomial degree at most
$\sum_i k_i\deg B_i$, just as in the polynomial case.

The convergence argument also extends. After replacing each finite source
mode by an independent marker, logarithmic composition uses
\[
 e^{i\omega(L+\log W)}=e^{i\omega L}W^{i\omega},
\]
which is analytic for $W$ near $1$ with its local logarithm. The implicit
derivative at the origin is still $p\ne0$. The finite-dimensional analytic
implicit function theorem gives an absolutely convergent marker expansion
for sufficiently small markers. Along the real source slice these markers
are bounded by constants times $z^{\delta_i}(1+|\log z|)^{d_i}$ and tend to
zero. The positive source gaps also make the forward derivative asymptotic
to $ap u^{p-1}$, so the selected real inverse branch exists near the endpoint.

The complete one-step outer boundary of the retained multi-index region
therefore bounds the omitted tail. Its least weight supplies the remainder
power, and the maximum summed polynomial degree over that finite boundary
supplies a conservative logarithmic envelope. This bound need not be sharp
after cancellation, but remains valid even at zeros of every displayed
oscillatory coefficient. A declared forward remainder is transported with
the same observable-dependent exponent as Theorem~\ref{thm:transport}; its
matching derivative hypothesis remains required. Constants and source
thresholds are not computed numerical certificates.

\subsection{Examples and scope}

For $f(x)=x+x^2\sin(\log x)$ at $0^+$, the inverse begins
\[
 y-y^2\sin L+y^3\sin L(2\sin L+\cos L)+\OO(y^4),
 \qquad L=\log y.
\]
With an additional $x^3\cos(2\log x)$, the coefficient of $y^3$ acquires
$-\cos(2L)$, combining a directly supplied frequency with frequencies
generated by convolution. The identity
\[
 \sin(\sqrt2L)\sin(\sqrt3L)
 +\tfrac12\cos\!\bigl(\sqrt{5+2\sqrt6}\,L\bigr)
 =\tfrac12\cos\!\bigl((\sqrt2-\sqrt3)L\bigr)
\]
exhibits an algebraic resonance and the resulting exact cancellation.

The cutoff is exclusive in the positive target power coordinate, with the
same finite-endpoint, infinite-endpoint, and observable conventions as in
the ordinary power--logarithmic case. Frequencies are fixed real constants;
parameter-dependent frequency collisions require their equality relations
to be resolved before forming complete blocks. Nonlinear logarithmic phases
require a different coefficient algebra.

The nonoscillatory envelope is essential even when a retained coefficient
vanishes along a sequence of source points. For example, the zeros of
$\sin(\log u)$ do not justify replacing a tail bound by an oscillatory
factor which also vanishes there. Bounds must control the entire omitted
sum uniformly near the endpoint, independently of such zeros.
